\chapter{Introducción}
El presente trabajo práctico se enfoca en el análisis del funcionamiento de diodos rectificadores y diodos Zener, considerando su comportamiento eléctrico bajo distintas condiciones de trabajo. A partir de esto, se busca comprender cómo responden estos dispositivos dentro de circuitos básicos y qué parámetros permiten caracterizar su operación.

Para abordar el estudio, se parte de la resolución analítica de los circuitos propuestos, obteniendo valores de tensión, corriente y potencia mediante modelos teóricos. Luego, dichos resultados se contrastan con simulaciones, lo que permite verificar la coherencia de los cálculos realizados y anticipar el comportamiento real del circuito.

Como etapa final, se implementan los circuitos en el laboratorio, realizando mediciones directas con instrumentos como el multímetro y el osciloscopio. Esta instancia experimental permite comparar los resultados obtenidos con los valores teóricos y simulados, identificando posibles desviaciones.

En conjunto, este trabajo práctico no solo permite estudiar el funcionamiento de los diodos en condiciones reales, sino también reforzar conceptos básicos de electrónica aplicados en cursos anteriores, necesarios para el desarrollo de contenidos más avanzados en la carrera.